% ==========================================================================
%  Chapter 7 — Applications and Use Cases
% ==========================================================================
\chapter{Applications and Use Cases}
\label{ch:applications}

\epigraph{%
  The value of UML is not that it virtualizes hardware---it's that it
  makes the kernel a first-class debugging target.%
}{Jeff Dike, Ottawa Linux Symposium, 2003}

User-Mode Linux's unique position in the virtualization design space---a
real kernel running as a debuggable userspace process---gives it a set
of application domains that no other single technology covers.  This
chapter surveys seven categories of use, ranging from established
practice (kernel debugging, network simulation) to emerging
opportunities made newly practical by the v2 redesign (coverage-guided
fuzzing, record/replay, snapshot-based iteration).

For each domain we describe the use case, explain why \UML{} is suited
to it, identify the limitations of current alternatives, and where
applicable, sketch the concrete workflow.


% ------------------------------------------------------------------
\section{Kernel Development and Debugging}
\label{sec:app-debugging}
% ------------------------------------------------------------------

The most natural use of \UML{} is the one Jeff Dike built it for:
running a Linux kernel under an ordinary debugger.

\subsection{The debugging workflow}

Because the \UML{} kernel is an ELF binary linked with full DWARF
debug information, a developer can attach \texttt{gdb} to it exactly
as one would attach to any other process:

\begin{lstlisting}[language=bash,caption={Launching UML under gdb}]
# Build with debug info
make ARCH=um O=build-um defconfig
make ARCH=um O=build-um -j$(nproc)

# Launch under gdb
gdb --args build-um/vmlinux mem=512M \
    root=/dev/root rootfstype=hostfs rootflags=/
\end{lstlisting}

From the \texttt{gdb} prompt the developer can:

\begin{itemize}[leftmargin=2em]
  \item Set breakpoints in any kernel function:
        \texttt{break do\_sys\_open}, \texttt{break \_\_schedule}.
  \item Single-step through the page fault handler, the scheduler,
        or the VFS layer.
  \item Inspect kernel data structures: \texttt{print *current},
        \texttt{print runqueues->cfs.nr\_running}.
  \item Use hardware watchpoints to catch unexpected writes to specific
        memory addresses.
  \item Examine the state of every ``CPU'' (host thread) using
        \texttt{info threads} and \texttt{thread apply all bt}.
\end{itemize}

This is qualitatively different from debugging a bare-metal or
QEMU-hosted kernel.  On bare metal, kernel debugging requires a serial
console, a JTAG probe, or KGDB---all of which add latency, limit the
set of available debugger commands, and require physical access to the
target machine.  With QEMU, one can use \texttt{gdb} over the GDB
remote protocol (\texttt{-s -S} flags), but the debugger operates on
the QEMU process, not on native host memory.  Variable inspection
requires the debugger to translate guest virtual addresses through
emulated page tables, and performance-sensitive operations like
conditional breakpoints incur orders-of-magnitude slowdowns due to
the remote protocol round-trips.

With \UML{}, the kernel \emph{is} the debuggee.  \texttt{gdb} reads
kernel memory through the same \texttt{ptrace} / \texttt{process\_vm\_readv}
mechanisms it uses for any local process.  Conditional breakpoints
evaluate at native speed.  The full repertoire of \texttt{gdb} Python
scripting is available for custom pretty-printers, breakpoint commands,
and automated inspection.

\subsection{kprobes and ftrace integration}

\UML{} supports the kernel's dynamic instrumentation frameworks.  A
developer can use \texttt{trace-cmd} from the host to record function
call graphs, schedule events, and lock contention traces.  Because
\UML{}'s \texttt{/sys/kernel/debug/tracing} is exposed through the
host filesystem (via \texttt{hostfs}), the developer can use familiar
host-side tools without installing anything inside the guest.

The kprobes subsystem allows inserting breakpoints at arbitrary kernel
addresses at runtime without recompilation.  Combined with eBPF~\cite{ebpf},
this enables sophisticated monitoring and analysis scenarios---for
example, counting the number of times a specific lock is contended per
second, or recording the arguments to every \syscall{open} call that
returns \texttt{EACCES}.

\subsection{Sanitizer support}

\UML{} supports KASAN~\cite{kasan} (Kernel Address Sanitizer),
KMSAN~\cite{kmsan} (Kernel Memory Sanitizer), and
KCSAN~\cite{kcsan} (Kernel Concurrency Sanitizer).  Because the
kernel runs as a userspace process, the shadow memory regions that
these sanitizers require can be allocated through ordinary
\syscall{mmap} calls without the complex early-boot memory management
that sanitizers require on hardware architectures.  This makes \UML{}
an ideal platform for sanitizer-heavy development workflows where
every memory access is checked for validity.

\begin{figure}[ht]
\centering
\begin{tikzpicture}[
    box/.style={
      draw=UMLGrid,
      fill=#1,
      minimum width=3.8cm,
      minimum height=1.0cm,
      rounded corners=3pt,
      font=\small,
      align=center
    },
    arrow/.style={
      -{Stealth[length=5pt]},
      thick,
      color=UMLMuted
    }
  ]

  % Developer tools
  \node[box=UMLBlue!8]  (gdb)    at (0,0)     {\textbf{gdb} / \textbf{rr}};
  \node[box=UMLTeal!8]  (trace)  at (4.5,0)   {\textbf{trace-cmd}\\ftrace};
  \node[box=UMLGreen!8] (san)    at (9,0)     {\textbf{KASAN}\\KMSAN / KCSAN};

  % UML kernel
  \node[box=KernelGold!8,minimum width=12.8cm] (uml) at (4.5,-1.8)
    {\textbf{UML Kernel} (single ELF binary with DWARF info)};

  % Arrows
  \draw[arrow] (gdb.south)   -- ++(0,-0.4) -| ($(uml.north west)+(1.5,0)$);
  \draw[arrow] (trace.south) -- (uml.north);
  \draw[arrow] (san.south)   -- ++(0,-0.4) -| ($(uml.north east)+(-1.5,0)$);

\end{tikzpicture}
\caption[Kernel debugging ecosystem under UML]{%
  Three complementary debugging approaches available when running the
  kernel under \UML{}: interactive debugging with \texttt{gdb} (or
  \texttt{rr} for record/replay), dynamic tracing with ftrace and
  \texttt{trace-cmd}, and compile-time sanitizers (KASAN, KMSAN, KCSAN).
  All three operate on the same kernel binary and can be used
  simultaneously.
}
\label{fig:app-debug-ecosystem}
\end{figure}


% ------------------------------------------------------------------
\section{Kernel Fuzzing}
\label{sec:app-fuzzing}
% ------------------------------------------------------------------

Coverage-guided kernel fuzzing is the most effective automated technique
for finding kernel bugs.  Google's syzkaller~\cite{syzkaller,vyukov2015}
is the world's most widely used kernel fuzzer, responsible for
discovering thousands of bugs in the Linux kernel since its introduction
in 2015.  Understanding how \UML{} fits into the fuzzing landscape
requires examining both the current state of the art and the
performance bottlenecks that limit it.

\subsection{The current landscape: QEMU-KVM}

Syzkaller's standard deployment targets QEMU/KVM virtual machines.
Each fuzzing instance boots a kernel inside a QEMU VM, runs a
\texttt{syz-executor} process that generates and executes syscall
programs, and collects coverage feedback via KCOV
(\texttt{CONFIG\_KCOV}).  The fuzzer mutates syscall sequences based
on which new code paths were exercised.

This architecture works well but imposes significant overhead:

\begin{enumerate}[leftmargin=2em]
  \item \textbf{Boot time.}  A QEMU/KVM instance takes 2--5 seconds to
        boot to the point where the fuzzer can begin executing syscalls.
  \item \textbf{Hardware requirement.}  QEMU/KVM requires
        \texttt{/dev/kvm}, which is unavailable in many CI/CD environments,
        nested virtualization scenarios, and shared cloud instances.
  \item \textbf{Communication overhead.}  The fuzzer communicates with
        the executor inside the VM through a virtio-serial channel or SSH.
        Each syscall program dispatch involves cross-VM communication.
  \item \textbf{Reset time.}  After a crash, the VM must be rebooted
        or restored from a snapshot, adding seconds of dead time.
\end{enumerate}

\subsection{The LKL precedent}

The Linux Kernel Library~\cite{lkl,purdila2010} demonstrated what
becomes possible when these overheads are eliminated.  By compiling the
kernel as a library and executing syscalls as function calls within the
same process, LKL-based fuzzers achieved approximately 60$\times$ the
throughput of QEMU-based approaches.  The LKL fuzzer can execute millions
of syscall sequences per second because there is no VM boundary, no boot
process, and no communication channel---the fuzzer and the kernel share
an address space.

However, LKL achieves this speed by sacrificing fidelity.  There is no
user/kernel separation, no process scheduling, no realistic signal
delivery, and no memory protection.  Bugs that depend on race conditions
between user and kernel execution, memory mapping semantics, or signal
handling are invisible to an LKL-based fuzzer.  The attack surface that
LKL exposes is a strict subset of the real kernel's attack surface.

\subsection{UML as a fuzzing target}

\UML{} occupies the sweet spot between QEMU and LKL:

\begin{table}[ht]
\centering
\caption{Fuzzing target comparison}
\label{tab:fuzz-comparison}
\begin{tabularx}{\textwidth}{l Y Y Y}
\toprule
\textbf{Property}  & \textbf{QEMU/KVM} & \textbf{LKL} & \textbf{UML (v2)} \\
\midrule
Boot time         & 2--5\,s           & N/A (library)    & $<$\,500\,ms \\
Syscall fidelity  & Full              & Partial          & Full \\
Process isolation & Full (HW)         & None             & Full (SW) \\
KCOV support      & Yes               & Custom           & Yes \\
Hardware requirement & /dev/kvm       & None             & None \\
Reset mechanism   & VM reboot/snapshot & Process restart & Snapshot/restore \\
Reset time        & 2--5\,s           & $<$\,1\,ms       & $<$\,1\,ms \\
\bottomrule
\end{tabularx}
\end{table}

With the v2 redesign's snapshot/restore capability, a \UML{} fuzzing
instance can be reset to a known-good state in under a millisecond.
The instance boots once, reaches the fuzzing entry point, takes a
snapshot, and then restores to that snapshot after each syscall sequence.
This eliminates boot time from the critical path entirely.

\subsection{The syzkaller vm/uml backend}

Syzkaller's architecture already supports pluggable VM backends through
its \texttt{vm/} package.  Adding a \UML{} backend requires implementing
the \texttt{vmimpl.Pool} and \texttt{vmimpl.Instance} interfaces:

\begin{itemize}[leftmargin=2em]
  \item \texttt{Create()} --- launch a \UML{} instance with the fuzzing
        kernel and root filesystem.
  \item \texttt{Forward()} --- set up port forwarding (trivial with \UML{}'s
        host networking).
  \item \texttt{Copy()} --- copy the \texttt{syz-executor} binary into the
        guest (via \texttt{hostfs}).
  \item \texttt{Run()} --- execute a command inside the guest and capture
        output.
  \item \texttt{Close()} --- terminate the instance.
\end{itemize}

Because \UML{} guest filesystems can be mounted via \texttt{hostfs}---a
direct pass-through of the host filesystem---there is no need for the
SSH or serial-port communication channels that QEMU requires.  The
fuzzer can write syscall programs directly to a shared directory and
read coverage data from the same path.

\subsection{KCOV integration}

KCOV (\texttt{CONFIG\_KCOV}) instruments the kernel to record which
basic blocks are executed during each syscall.  \UML{} supports KCOV
natively because KCOV is architecture-independent---it hooks into the
compiler's instrumentation callbacks, not into architecture-specific
trap handling.  The \texttt{/sys/kernel/debug/kcov} interface is
available inside the \UML{} guest and can be accessed through
\texttt{hostfs} from the host-side fuzzer.


% ------------------------------------------------------------------
\section{Network Simulation}
\label{sec:app-network}
% ------------------------------------------------------------------

One of the most successful real-world deployments of \UML{} has been
in network simulation and education, where the ability to run many
independent Linux instances---each with its own full networking
stack---without any hardware virtualization support has proven
invaluable.

\subsection{Netkit}

Netkit~\cite{pizzonia2008}, developed by Maurizio Pizzonia and Massimo
Rimondini at Roma Tre University, is the most widely used \UML{}-based
networking tool.  Each Netkit ``virtual machine'' is a \UML{} instance
connected to virtual Ethernet segments implemented as UNIX domain sockets.
A lab configuration file specifies the topology: which machines exist,
which interfaces they have, and which virtual collision domains they
connect to.

A typical Netkit lab might define a topology with:

\begin{itemize}[leftmargin=2em]
  \item Three routers running OSPF (via Quagga or FRRouting).
  \item Two Ethernet segments with multiple hosts.
  \item A NAT gateway connecting the virtual topology to the host network.
  \item A DNS server and a web server on separate subnets.
\end{itemize}

Each of these ``machines'' is a \UML{} process consuming 32--64\,MB of
memory.  A laptop can run dozens of instances simultaneously, enabling
students to build and experiment with network topologies that would
otherwise require a rack of physical hardware.

Netkit has been adopted by universities in Italy, Brazil, Germany,
Australia, and elsewhere for teaching courses on computer networking,
network security, and system administration.  The Netkit-NG fork
updated the project for modern kernels, and Kathara (a Docker-based
reimagining) continues the tradition, though with a shift away from
\UML{} toward containers.

\subsection{Marionnet}

Marionnet~\cite{marionnet}, developed at Paris 13 University, provides
a graphical interface for building virtual network topologies using
\UML{}.  Students drag and drop routers, switches, and hosts onto a
canvas, connect them with virtual cables, and then start the entire
topology.  Each node runs a full Linux instance under \UML{}, providing
realistic behavior including ARP, ICMP, TCP connection establishment,
and routing protocol convergence.

\subsection{Wireless network simulation}

Intel's WiFi team has used \UML{} for wireless network simulation,
running multiple \UML{} instances with virtual wireless interfaces
(mac80211\_hwsim) to test WiFi drivers and the \texttt{cfg80211}/\texttt{mac80211}
stack without physical wireless hardware.  This allows automated testing
of scenarios like roaming, channel switching, power management, and
multi-AP coordination that would be difficult to reproduce reliably
with physical hardware.

\begin{figure}[ht]
\centering
\begin{tikzpicture}[
    host/.style={
      draw=UMLGrid,
      fill=UMLGreen!8,
      minimum width=2.0cm,
      minimum height=0.9cm,
      rounded corners=3pt,
      font=\small\bfseries
    },
    switch/.style={
      draw=UMLGrid,
      fill=UMLTeal!8,
      minimum width=2.0cm,
      minimum height=0.7cm,
      rounded corners=3pt,
      font=\small
    },
    router/.style={
      draw=UMLGrid,
      fill=UMLBlue!8,
      minimum width=2.0cm,
      minimum height=0.9cm,
      rounded corners=3pt,
      font=\small\bfseries
    },
    link/.style={
      thick,
      color=UMLMuted
    }
  ]

  % Router
  \node[router] (r1) at (4.5,2.5) {Router R1};

  % Switches
  \node[switch] (sw1) at (1.5,1.2) {Switch SW1};
  \node[switch] (sw2) at (7.5,1.2) {Switch SW2};

  % Hosts
  \node[host] (h1) at (0,0) {Host A};
  \node[host] (h2) at (3,0) {Host B};
  \node[host] (h3) at (6,0) {Host C};
  \node[host] (h4) at (9,0) {Host D};

  % Links
  \draw[link] (r1) -- (sw1);
  \draw[link] (r1) -- (sw2);
  \draw[link] (sw1) -- (h1);
  \draw[link] (sw1) -- (h2);
  \draw[link] (sw2) -- (h3);
  \draw[link] (sw2) -- (h4);

  % Labels
  \node[font=\scriptsize\itshape,color=UMLMuted,below=0.1cm of h1]
    {UML instance};
  \node[font=\scriptsize\itshape,color=UMLMuted,below=0.1cm of h4]
    {UML instance};
  \node[font=\scriptsize\itshape,color=UMLMuted,above=0.1cm of r1]
    {UML instance (OSPF)};

\end{tikzpicture}
\caption[Netkit-style network topology]{%
  A simple Netkit-style network topology.  Each host, switch, and
  router is an independent \UML{} instance.  Virtual Ethernet
  segments are implemented as UNIX domain sockets on the host.
  The entire topology runs as ordinary userspace processes without
  any hardware virtualization.
}
\label{fig:app-netkit}
\end{figure}


% ------------------------------------------------------------------
\section{Security Research}
\label{sec:app-security}
% ------------------------------------------------------------------

Security researchers need environments where they can safely execute
kernel exploits, analyze malware behavior at the kernel level, and
develop mitigations---all without risking the host system.  \UML{}
provides this environment with properties that neither containers nor
traditional VMs fully match.

\subsection{Kernel exploit development and analysis}

When analyzing a kernel vulnerability---say, a use-after-free in a
filesystem driver---a researcher needs to:

\begin{enumerate}[leftmargin=2em]
  \item Build a kernel with the vulnerable code and full debug symbols.
  \item Trigger the vulnerability reliably.
  \item Observe the kernel state at the moment of corruption.
  \item Develop and test a mitigation.
\end{enumerate}

With \UML{}, every step happens in a single terminal session.  The
researcher builds the vulnerable kernel as an ELF binary, runs it
under \texttt{gdb}, sets a watchpoint on the freed memory region,
and observes the exact instruction that performs the use-after-free.
The kernel's full DWARF debug information is available, so the
researcher can inspect the call stack, the state of the slab
allocator, and the contents of every relevant data structure.

\subsection{Advantage over containers}

Containers (Docker, Kubernetes) share the host kernel.  A kernel
exploit that works in a container works against the host.  This is
not a theoretical concern: container escape vulnerabilities
(CVE-2019-5736, CVE-2020-15257, CVE-2022-0185) regularly demonstrate
that namespace and cgroup isolation do not constitute a security
boundary at the kernel level.

\UML{} provides genuine kernel isolation.  Each \UML{} instance
runs its own kernel in a separate address space.  A kernel exploit
that achieves code execution inside a \UML{} instance has compromised
a userspace process from the host's perspective.  It has no access
to the host kernel's data structures, no ability to manipulate host
page tables, and no path to privilege escalation on the host.  The
host kernel's normal process isolation guarantees contain the
compromised \UML{} instance.

\subsection{Honeypot deployment}

\UML{}'s lightweight footprint makes it practical to deploy
kernel-level honeypots.  Each honeypot is a \UML{} instance running
a deliberately vulnerable kernel, configured to log all system calls
and network traffic.  Because \UML{} instances are ordinary processes,
they can be monitored, resource-limited (via cgroups), and terminated
from the host without any special hypervisor tooling.  The host can
run dozens or hundreds of honeypot instances, each presenting a
different kernel version or configuration to probe for attacker
behavior.

\subsection{Malware analysis}

Kernel-level malware (rootkits, bootkits) requires a real kernel
environment for analysis.  Running such malware in a container is
dangerous because the malware has direct access to the host kernel.
Running it in a QEMU VM is safe but makes observation difficult:
the analyst must use the GDB remote protocol or custom QEMU plugins
to inspect kernel state.  \UML{} provides the best of both
worlds---a real kernel that the malware can infect, with the full
observability of a native userspace process.


% ------------------------------------------------------------------
\section{CI/CD and Automated Testing}
\label{sec:app-cicd}
% ------------------------------------------------------------------

Continuous integration and testing of kernel code requires an
environment that can boot a kernel, run tests, and report results.
The dominant approach uses QEMU/KVM, but this has a fundamental
limitation: it requires access to \texttt{/dev/kvm}.

\subsection{The /dev/kvm problem}

Many CI/CD environments lack hardware virtualization support:

\begin{description}[style=nextline,leftmargin=2em]
  \item[GitHub Actions (default runners).]
    The standard GitHub-hosted runners do not expose \texttt{/dev/kvm}.
    Larger runners with KVM access are available at increased cost.
  \item[Shared cloud instances.]
    Many cloud providers do not support nested virtualization on their
    standard instance types.  Even when nested virtualization is
    available, it incurs significant performance overhead.
  \item[Containerized CI.]
    CI systems that run jobs inside containers (Jenkins with Docker,
    GitLab CI) typically do not pass through \texttt{/dev/kvm} to
    the container for security reasons.
  \item[Locked-down workstations.]
    Corporate environments may disable hardware virtualization in BIOS
    or restrict access to \texttt{/dev/kvm} via udev rules.
\end{description}

\UML{} with the \Seccomp{} or \texttt{ptrace} backend has no such
requirement.  It runs wherever a Linux userspace process can run.
This makes it uniquely suited for kernel testing in environments where
hardware-assisted virtualization is unavailable.

\subsection{kselftest integration}

The kernel's self-test framework (\texttt{tools/testing/selftests/})
includes tests that require a running kernel environment.  The v2
redesign adds a comprehensive test suite under
\umlpath{tools/testing/selftests/um/} that exercises backend lifecycle
management, syscall interception, SMP behavior, and memory management.

These tests can be run as part of a standard CI pipeline:

\begin{lstlisting}[language=bash,caption={Running kselftest under UML in CI}]
# Build UML kernel with kselftest support
make ARCH=um O=build-um defconfig
make ARCH=um O=build-um -j$(nproc)

# Run the UML-specific selftests
make ARCH=um O=build-um -C tools/testing/selftests/um run_tests

# Run broader kernel selftests under UML
./build-um/vmlinux mem=512M \
    init=/path/to/run-selftests.sh
\end{lstlisting}

\subsection{Linux Test Project (LTP)}

The Linux Test Project (LTP) is a comprehensive test suite covering
syscall behavior, filesystem semantics, memory management, IPC, and
networking.  LTP tests can run inside a \UML{} instance with no
modifications, exercising the real kernel code paths.  Because \UML{}
boots in under a second, a CI pipeline can boot a fresh \UML{} instance
for each test group, providing isolation without the multi-second
overhead of QEMU boot cycles.

\subsection{Bisection workflows}

When a kernel regression is identified, \texttt{git bisect} requires
building and testing many kernel versions.  With QEMU, each test
iteration takes 10--30 seconds (build + boot + test + shutdown).
With \UML{}, the same cycle takes 2--5 seconds because:

\begin{enumerate}[leftmargin=2em]
  \item \UML{} builds are faster (smaller architecture-specific code).
  \item Boot time is under 500\,ms.
  \item No VM teardown is needed---the process simply exits.
\end{enumerate}

Over a typical bisection of 12--15 steps, this difference compounds
from minutes with QEMU to under a minute with \UML{}.


% ------------------------------------------------------------------
\section{Education}
\label{sec:app-education}
% ------------------------------------------------------------------

Operating systems courses have a persistent pedagogical problem: the
gap between theory and practice.  Students learn about schedulers,
virtual memory, and file systems in lecture, but interacting with a
real kernel traditionally required dedicated hardware, serial console
access, and the risk of bricking a machine.

\subsection{UML as a teaching tool}

\UML{} eliminates this barrier.  Every student can:

\begin{enumerate}[leftmargin=2em]
  \item Clone the kernel source tree on their laptop.
  \item Build it as a \texttt{ARCH=um} binary in minutes
        (faster than a full x86 build because the architecture-specific
        code is smaller).
  \item Run the kernel in a terminal window.
  \item Modify kernel code, rebuild, and re-run---a cycle measured in
        seconds.
  \item Attach \texttt{gdb} to observe the scheduler making decisions,
        the page fault handler allocating pages, or the VFS dispatching
        I/O to a filesystem.
\end{enumerate}

There is no risk to the student's laptop.  The worst that can happen
is the \UML{} process crashes---which is itself an educational
experience, as the student can examine the core dump.

\subsection{Concrete course applications}

\begin{description}[style=nextline,leftmargin=2em]
  \item[Process scheduling.]
    Students can add \texttt{printk} statements to the CFS scheduler,
    rebuild in seconds, and observe the scheduling decisions made for
    competing workloads.  With \texttt{gdb}, they can single-step
    through \texttt{pick\_next\_task\_fair()} and watch the red-black
    tree traversal in real time.

  \item[Virtual memory.]
    Students can trace page fault handling from the initial trap through
    \texttt{handle\_mm\_fault()}, page table walk, page allocation, and
    PTE installation.  \UML{}'s memory management uses host \syscall{mmap}
    calls, making it possible to correlate kernel-level page faults with
    host-level memory operations.

  \item[Filesystem implementation.]
    Students can build a minimal filesystem module, load it into a
    running \UML{} instance, and trace the complete path of a
    \syscall{read} call from the system call entry point through the
    VFS layer, the filesystem-specific \texttt{read\_iter} implementation,
    and the block layer.

  \item[Device drivers.]
    \UML{}'s virtual device model allows students to write simple
    character and block device drivers that interact with host-side
    files.  The driver runs in the \UML{} kernel, and its behavior
    can be observed from both the guest perspective (via
    \texttt{/dev/} entries) and the host perspective (via
    \texttt{gdb} and \texttt{strace}).

  \item[Concurrency and synchronization.]
    With SMP support in the v2 redesign, students can observe real
    lock contention, priority inversion, and RCU grace periods.  KCSAN
    can be enabled to catch data races, turning abstract concurrency
    bugs into concrete, reproducible reports.
\end{description}

\subsection{Comparison with teaching alternatives}

\begin{table}[ht]
\centering
\caption{Comparison of OS teaching environments}
\label{tab:edu-comparison}
\begin{tabularx}{\textwidth}{l Y Y Y Y}
\toprule
\textbf{Property} & \textbf{xv6} & \textbf{QEMU/KVM} & \textbf{Minix} & \textbf{UML} \\
\midrule
Real kernel code    & No   & Yes  & No   & Yes \\
Native debugging    & No   & GDB remote & No & GDB native \\
Rebuild cycle       & Fast & Slow & Medium & Fast \\
Hardware required   & None & /dev/kvm & None & None \\
Industry relevance  & Low  & High & Low  & High \\
\bottomrule
\end{tabularx}
\end{table}

xv6 (MIT's teaching OS) is elegant and well-documented, but it is
not Linux.  Skills learned in xv6 do not directly transfer to
kernel development.  QEMU/KVM runs real Linux but requires hardware
virtualization and imposes slow iteration cycles.  \UML{} is the
only option that combines the real Linux kernel with the fast
iteration and native debuggability that an educational environment
demands.


% ------------------------------------------------------------------
\section{Record/Replay for Bug Reproduction}
\label{sec:app-record-replay}
% ------------------------------------------------------------------

Many of the most difficult kernel bugs are non-deterministic: they
depend on the precise timing of interrupts, the order of concurrent
operations, or the exact sequence of I/O completions.  Reproducing
such bugs is often the hardest part of fixing them.  Record/replay
systems address this by capturing the non-deterministic inputs to a
computation during a ``record'' run and replaying them exactly during
a ``replay'' run, guaranteeing deterministic reproduction.

\subsection{Existing approaches}

\textbf{rr}~\cite{rr-project} (Mozilla) is the most successful
userspace record/replay system.  It records the results of system
calls, signal deliveries, and non-deterministic instructions
(\texttt{RDTSC}, \texttt{RDRAND}) during a recording session, and
replays them faithfully.  However, rr operates at the userspace level:
it can record and replay a userspace program, but it cannot record and
replay kernel execution.  If a bug is in the kernel, rr cannot help.

\textbf{ReVirt}~\cite{dunlap2002} demonstrated virtual machine-level
record/replay by logging all non-deterministic events (interrupts,
DMA, device I/O) entering a VM.  But ReVirt required modifications
to the hypervisor and imposed significant overhead.

\subsection{UML record/replay}

\UML{}'s architecture makes record/replay uniquely tractable.  The
sources of non-determinism in a \UML{} instance are precisely
enumerable:

\begin{enumerate}[leftmargin=2em]
  \item \textbf{Timer interrupts.}  Generated by host-side
        \texttt{SIGALRM} signals.  The record layer logs the exact
        instruction pointer at which each signal was delivered.
  \item \textbf{I/O completions.}  Disk reads, network packet arrivals,
        and other I/O events arrive as host system call return values.
        The record layer logs these values.
  \item \textbf{Inter-processor interrupts.}  In SMP mode, IPIs from
        other virtual CPUs introduce ordering non-determinism.  The
        record layer logs the IPI delivery order.
  \item \textbf{Host entropy.}  Calls to \syscall{getrandom} or reads
        from \texttt{/dev/urandom} on the host.  The record layer logs
        the returned bytes.
\end{enumerate}

During replay, the record layer intercepts each of these sources and
injects the logged values at the exact same points in execution.  The
result is bit-for-bit deterministic replay of the kernel's execution,
including all scheduler decisions, lock acquisitions, and memory
allocations.

\subsection{Time-travel debugging}

When combined with \texttt{gdb}, record/replay enables \emph{time-travel
debugging}: the ability to run the kernel execution backward.  The
developer can set a watchpoint on a corrupted data structure, run
backward to the point of corruption, and identify the exact code path
that caused it.  This is the kernel-level equivalent of rr's
\texttt{reverse-continue} command, applied to the kernel itself.

Jeff Dike implemented an early version of time-travel support for
\UML{}, and the v2 redesign extends this with SMP-aware recording
and the ability to replay across snapshot/restore boundaries.


% ------------------------------------------------------------------
\section{Snapshot/Restore for Fast Iteration}
\label{sec:app-snapshot}
% ------------------------------------------------------------------

The ability to capture the complete state of a running kernel and
restore it later is valuable for several workflows:

\subsection{Fuzzing reset}

As discussed in Section~\ref{sec:app-fuzzing}, coverage-guided fuzzers
need to reset the kernel to a clean state after each test iteration.
With snapshot/restore, this reset is a memory operation rather than a
reboot.  The v2 implementation captures:

\begin{itemize}[leftmargin=2em]
  \item All guest physical memory pages (via \syscall{mmap} snapshots).
  \item CPU register state for all virtual CPUs.
  \item Kernel timer state.
  \item Virtual device state (block device positions, network queue
        contents).
\end{itemize}

Restore replays these snapshots by remapping memory regions and
restoring register state.  Because \UML{}'s memory is ordinary
process memory (anonymous \syscall{mmap} regions), the snapshot can
use copy-on-write semantics: the restore operation marks all pages
as CoW, and only pages that are actually modified during the next
test iteration incur a copy.  This makes restore effectively
sub-microsecond for the common case where only a small fraction of
pages are modified between iterations.

\subsection{Debugging iteration}

A developer investigating a complex kernel bug can take a snapshot
just before the bug manifests, then repeatedly restore and re-examine
the bug with different breakpoints, watchpoints, or instrumentation.
This is similar to recording and replaying, but more flexible: the
developer can modify kernel code between restore iterations, recompile
individual modules, and reload them into the restored kernel.

\subsection{Test parallelization}

A single snapshot can be restored into multiple independent \UML{}
instances, each running a different test.  This enables embarrassingly
parallel testing: boot the kernel once, take a snapshot at the
test-ready state, and fork it into $N$ copies, each running one test
suite.  The boot cost is amortized across all test suites.

\begin{figure}[ht]
\centering
\begin{tikzpicture}[
    state/.style={
      draw=UMLGrid,
      fill=#1,
      minimum width=2.5cm,
      minimum height=0.8cm,
      rounded corners=3pt,
      font=\small
    },
    arrow/.style={
      -{Stealth[length=5pt]},
      thick,
      color=UMLMuted
    },
    darrow/.style={
      -{Stealth[length=5pt]},
      thick,
      dashed,
      color=UMLAmber
    }
  ]

  % Boot sequence
  \node[state=UMLBlue!8]  (boot)  at (0,0)     {Boot};
  \node[state=UMLGreen!8] (snap)  at (3,0)     {Snapshot};

  % Fork into multiple instances
  \node[state=KernelGold!8] (t1) at (7,1)   {Test Suite A};
  \node[state=KernelGold!8] (t2) at (7,0)   {Test Suite B};
  \node[state=KernelGold!8] (t3) at (7,-1)  {Test Suite C};

  % Restore arrows
  \draw[arrow] (boot) -- (snap);
  \draw[darrow] (snap) -- (t1) node[midway,above,font=\scriptsize] {restore};
  \draw[darrow] (snap) -- (t2);
  \draw[darrow] (snap) -- (t3) node[midway,below,font=\scriptsize] {restore};

  % Annotation
  \node[font=\scriptsize\itshape,color=UMLMuted] at (3,-1.5)
    {Boot cost paid once};
  \node[font=\scriptsize\itshape,color=UMLMuted] at (7,-1.9)
    {Each restore: $<$\,1\,$\mu$s (CoW)};

\end{tikzpicture}
\caption[Snapshot-based test parallelization]{%
  Snapshot-based test parallelization.  The kernel boots once and
  reaches a test-ready state.  A snapshot is taken and restored into
  multiple independent instances, each running a different test suite.
  Restore uses copy-on-write semantics for sub-microsecond latency.
}
\label{fig:app-snapshot-parallel}
\end{figure}


% ------------------------------------------------------------------
\section{Summary}
\label{sec:app-summary}
% ------------------------------------------------------------------

The applications surveyed in this chapter share a common theme:
\UML{} provides capabilities that require the \emph{combination} of
kernel fidelity, process-level debuggability, and minimal host
requirements---a combination that no other single technology offers.

\begin{table}[ht]
\centering
\caption{Application domains and key UML advantages}
\label{tab:app-summary}
\begin{tabularx}{\textwidth}{l Y Y}
\toprule
\textbf{Application} & \textbf{Key UML advantage} & \textbf{Alternative limitations} \\
\midrule
Kernel debugging   & Native gdb, full DWARF        & QEMU: remote protocol only \\
Kernel fuzzing     & No /dev/kvm, fast reset        & QEMU: slow boot; LKL: no isolation \\
Network simulation & Many instances, no HW virt     & Physical hardware: expensive \\
Security research  & Full kernel isolation           & Containers: shared kernel \\
CI/CD testing      & Runs anywhere Linux runs        & QEMU: requires /dev/kvm \\
Education          & Fast rebuild, native debug      & xv6: not real Linux \\
Record/replay      & Enumerable non-determinism      & rr: userspace only \\
Snapshot/restore   & Sub-$\mu$s CoW restore          & QEMU: multi-second restore \\
\bottomrule
\end{tabularx}
\end{table}

The v2 redesign strengthens \UML{}'s position in each of these domains.
The backend-ops abstraction allows choosing the optimal backend for each
use case (seccomp for portability, kvm-v2 for performance).  SMP support
enables testing of concurrent code paths.  Snapshot/restore and
record/replay transform \UML{} from a development convenience into a
serious research infrastructure.

Chapter~\ref{ch:related} examines how \UML{} compares to the related
systems that serve some of these same domains.
